\documentclass[hyperref=unicode, presentation,10pt]{beamer}

\usepackage[absolute,overlay]{textpos}
\usepackage{array}
\usepackage{graphicx}
\usepackage{adjustbox}
\usepackage{mhchem}
\usepackage{chemfig}
\usepackage{caption}

%dělení slov
\usepackage{ragged2e}
\let\raggedright=\RaggedRight
%konec dělení slov

\addtobeamertemplate{frametitle}{
	\let\insertframetitle\insertsectionhead}{}
\addtobeamertemplate{frametitle}{
	\let\insertframesubtitle\insertsubsectionhead}{}

\makeatletter
\CheckCommand*\beamer@checkframetitle{\@ifnextchar\bgroup\beamer@inlineframetitle{}}
\renewcommand*\beamer@checkframetitle{\global\let\beamer@frametitle\relax\@ifnextchar\bgroup\beamer@inlineframetitle{}}
\makeatother
\setbeamercolor{section in toc}{fg=red}
\setbeamertemplate{section in toc shaded}[default][100]

\usepackage{fontspec}
\usepackage{unicode-math}

\usepackage{polyglossia}
\setdefaultlanguage{czech}

\def\uv#1{„#1“}

\mode<presentation>{\usetheme{default}}
\usecolortheme{crane}

\setbeamertemplate{footline}[frame number]

\title[Crisis]
{C2062 -- Anorganická chemie II}

\subtitle{Demonstrační pokusy -- d-blok -- 3.--7. skupina}
\author{Zdeněk Moravec, hugo@chemi.muni.cz \\ \adjincludegraphics[height=60mm]{../IUPAC_PSP.jpg}}
\date{}

\begin{document}
\begin{frame}
	\titlepage
\end{frame}

\section{Demonstrační pokusy -- d-blok -- 3.--7. skupina}

\frame{
	\frametitle{}
	\vfill
	\textbf{Redukce chromanu vodíkem}
	\begin{itemize}
		\item \ce{2 CrO$_4^{2-}$ + 2 H^{+} -> Cr2O$_7^{2-}$ + H2O}
		\item \ce{Cr2O$_7^{2-}$ + 3 H2 + 8 H+ -> 2 Cr^{3+} + 7 H2O}
	\end{itemize}
	\vspace{5mm}
	\hrule
	\vspace{5mm}

	\textbf{Srážení chromanů}
	\begin{itemize}
		\item \ce{2 Ag+ + CrO$_4^{2-}$ -> Ag2CrO4}
		\item \ce{Pb^{2+} + CrO$_4^{2-}$ -> PbCrO4} (chromová žluť)
		\item \ce{2 Al^{3+} + 3 CrO$_4^{2-}$ -> Al2(CrO4)3}
	\end{itemize}
	\vspace{5mm}
	\hrule
	\vspace{5mm}
	\textbf{Komplex oxido-diperoxidochromový}
	\begin{itemize}
		\item \ce{Cr2O$_7^{2-}$ + 4 H2O2 + 2 H^{+} -> 2 [Cr(O)(O2)2] + 5 H2O}\footnote[frame]{\href{https://cs.wikibooks.org/wiki/Chemick\%C3\%A9_pokusy/Peroxokomplex_chromu}{Peroxokomplex chromu}}
	\end{itemize}
	\vfill
}

\frame{
	\frametitle{}
	\vfill
	\textbf{Reverzibilní přechod chroman--dichroman}
	\begin{itemize}
		\item \ce{2 CrO$_4^{2-}$ + 2 H^{+} -> Cr2O$_7^{2-}$ + H2O}
		\item \ce{Cr2O$_7^{2-}$ + OH^{-} -> 2 CrO$_4^{2-}$ + H+}
	\end{itemize}

	\begin{columns}
		\begin{column}{.6\textwidth}
			\begin{figure}
				\adjincludegraphics[width=\textwidth]{img/chroman-dichroman.png}
				\caption*{Struktura dichromanu a chromanu}
			\end{figure}
		\end{column}

		\begin{column}{.4\textwidth}
			\begin{figure}
				\adjincludegraphics[height=.35\textheight]{img/Chromate_dichromate_equilibrium.png}
				\caption*{Roztoky dichromanu a chromanu.\footnote[frame]{\href{https://commons.wikimedia.org/wiki/File:Chromate_dichromate_equilibrium.png}{Zdroj: Capaccio/Commons}}}
			\end{figure}
		\end{column}
	\end{columns}
	\vfill
}

\frame{
	\frametitle{}
	\vfill
	\begin{itemize}
		\item Dichroman amonný podléhá samovolnému rozkladu za uvolnění plynného dusíku a vzniku oxidu chromitého:\footnote[frame]{\href{https://commons.wikimedia.org/wiki/File:Thermal_decomposition_of_ammonium_dichromate_(1).jpg}{Zdroj: Rando Tuvikene/Commons}}$^,$\footnote[frame]{\href{https://commons.wikimedia.org/wiki/File:Thermal_decomposition_of_ammonium_dichromate_(2).jpg}{Zdroj: Rando Tuvikene/Commons}}$^,$\footnote[frame]{\href{https://commons.wikimedia.org/wiki/File:Thermal_decomposition_of_ammonium_dichromate_(3).jpg}{Zdroj: Rando Tuvikene/Commons}}
		\item \ce{(NH4)2Cr2O7 ->[200 $^\circ$C] Cr2O3 + N2 ^ + 4 H2O ^}
	\end{itemize}
	\begin{columns}
		\begin{column}{.33\textwidth}
			\begin{figure}
				\adjincludegraphics[width=.9\textwidth]{img/Sopka1.jpg}
			\end{figure}
		\end{column}

		\begin{column}{.33\textwidth}
			\begin{figure}
				\adjincludegraphics[width=.9\textwidth]{img/Sopka2.jpg}
			\end{figure}
		\end{column}

		\begin{column}{.33\textwidth}
			\begin{figure}
				\adjincludegraphics[width=.9\textwidth]{img/Sopka3.jpg}
			\end{figure}
		\end{column}
	\end{columns}
	\vfill
}

\frame{
	\frametitle{}
	\vfill
	\textbf{Alkalické tavení \ce{MnO2}}

	\begin{itemize}
		\item Oxid manganičitý (\ce{MnO2}) je možné oxidovat kyslíkem v bazickém prostředí na manganan.
		\item \ce{2 KNO3 -> 2 KNO2 + O2}
		\item \ce{2 MnO2 + 4 KOH + O2 -> 2 K2MnO4 + 2 H2O}
		\item Ten následně disproporcionuje za vzniku manganistanu a oxidu manganičitého.
		\item \ce{3 K2MnO4 + 4 HCl -> 2 KMnO4 + MnO2 + 4 KCl + 2 H2O}
	\end{itemize}
	\vfill
}

\end{document}